\documentclass[12pt,a4paper]{article}

% Paquetes
\usepackage[utf8]{inputenc}
\usepackage[spanish]{babel}
\usepackage{geometry}
\geometry{margin=2.5cm}
\usepackage{hyperref}
\usepackage{xcolor}
\usepackage{graphicx}
\usepackage{fancyhdr}
\usepackage{titlesec}
\usepackage{enumitem}
\usepackage{booktabs}
\usepackage{array}

% Colores personalizados
\definecolor{azulgit}{RGB}{24,80,50}
\definecolor{grisclaro}{RGB}{245,245,245}

% Configuración de hipervínculos
\hypersetup{
    colorlinks=true,
    linkcolor=azulgit,
    urlcolor=azulgit,
    citecolor=azulgit
}

% Encabezado y pie de página
\pagestyle{fancy}
\fancyhf{}
\fancyhead[L]{\textcolor{azulgit}{\small Programadores Por La Paz}}
\fancyhead[R]{\textcolor{azulgit}{\small Semana 1}}
\fancyfoot[C]{\thepage}
\renewcommand{\headrulewidth}{0.4pt}

% Títulos
\titleformat{\section}
  {\normalfont\Large\bfseries\color{azulgit}}
  {\thesection}{1em}{}
\titleformat{\subsection}
  {\normalfont\large\bfseries\color{azulgit!80!black}}
  {\thesubsection}{1em}{}

\begin{document}

% Título
\begin{titlepage}
    \centering
    \vspace*{2cm}
    
    {\Huge\bfseries\textcolor{azulgit}{DOCUMENTO DE PRESENTACIÓN}\\[0.5cm]}
    {\Large\textcolor{azulgit!70!black}{Tarea Semana 1}}\\[1cm]
    
    \vspace{1cm}
    
    {\Large Ciudadanía Digital y Herramientas Técnicas}\\[2cm]
    
    \vspace{2cm}
    
    {\large\textbf{Programa:}}\\
    {\large Programadores Por La Paz}\\[1.5cm]
    
    {\large\textbf{Universidad de Cartagena}}\\[2cm]
    
    \vfill
    
    {\large\textbf{Estudiante:}}\\
    {\Large Devin Polanco Romero}\\[1cm]
    
    {\large Junio 2026}
    
\end{titlepage}

% Índice
\tableofcontents
\newpage

% Contenido principal
\section{Información General}

\begin{table}[h]
\centering
\begin{tabular}{|>{\bfseries}l|l|}
\hline
\textbf{Campo} & \textbf{Detalle} \\
\hline
Estudiante & Devin Polanco Romero \\
\hline
Programa & Programadores Por La Paz \\
\hline
Institución & Universidad de Cartagena \\
\hline
Actividad & Tarea Semana 1 \\
\hline
Tema & Ciudadanía Digital y Herramientas Técnicas \\
\hline
Fecha & Junio 2026 \\
\hline
\end{tabular}
\caption{Información de la actividad}
\end{table}

\section{Enlace del Repositorio}

El desarrollo de la actividad fue publicado en el siguiente repositorio de GitHub:

\vspace{0.5cm}

\begin{center}
\large
\url{https://github.com/DevCorpSoftwareFactory/Programadores_Para_La_Paz_Semana_1}
\end{center}

\vspace{0.5cm}

\noindent La carpeta \texttt{semana1} contiene todos los archivos solicitados para la entrega de la actividad.

\section{Objetivo de la Semana}

Comprender la relación entre \textbf{ciudadanía digital}, uso responsable de la información, \textbf{pensamiento algorítmico} y organización técnica del trabajo, iniciando el uso de Git y GitHub para registrar y publicar la actividad realizada.

\section{Desarrollo de la Actividad}

\subsection{Preguntas de Selección Múltiple}

Se respondieron tres preguntas sobre ciudadanía digital y pensamiento algorítmico:

\begin{enumerate}[label=\textbf{\arabic*.}]
    \item \textbf{¿Qué es ciudadanía digital?} \\
    Respuesta: El uso responsable de la tecnología para informarse, participar y compartir información de manera ética.
    
    \item \textbf{¿Por qué es importante verificar la información?} \\
    Respuesta: Para evitar difundir información falsa o engañosa.
    
    \item \textbf{¿Qué es un algoritmo?} \\
    Respuesta: Una lista de pasos ordenados para resolver un problema.
\end{enumerate}

\subsection{Algoritmo: Compartir Información Confiable}

Se diseñó un algoritmo de 8 pasos para compartir información responsable en internet:

\begin{enumerate}
    \item Identificar la información y definir el tema
    \item Buscar en fuentes oficiales y reconocidas
    \item Verificar las credenciales del autor
    \item Contrastar con al menos 2 fuentes adicionales
    \item Revisar la fecha de publicación
    \item Evaluar el lenguaje (descartar sensacionalismo)
    \item Verificar datos verificables y referencias
    \item Compartir solo cuando se confirme su veracidad
\end{enumerate}

\subsection{Reflexión}

Se desarrolló una reflexión sobre la importancia de la ciudadanía digital, destacando:

\begin{itemize}
    \item La responsabilidad de verificar la información antes de compartirla
    \item La conexión entre participación ciudadana y entornos digitales confiables
    \item El papel de Git en la organización del trabajo técnico
    \item La importancia de la ética y el compromiso con la verdad
\end{itemize}

\section{Archivos Entregados}

La carpeta \texttt{semana1} contiene los siguientes archivos:

\begin{table}[h]
\centering
\begin{tabular}{|l|l|}
\hline
\textbf{Archivo} & \textbf{Contenido} \\
\hline
preguntas-semana1.txt & Respuestas de selección múltiple \\
\hline
algoritmo-semana1.txt & Algoritmo de 8 pasos para compartir información \\
\hline
reflexion-semana1.txt & Reflexión sobre ciudadanía digital \\
\hline
\end{tabular}
\caption{Archivos de la actividad}
\end{table}

\section{Conclusiones}

\begin{itemize}
    \item Se comprendió el concepto de ciudadanía digital y su importancia en el entorno actual.
    \item Se diseñó un algoritmo práctico para verificar y compartir información confiable.
    \item Se reconoció la conexión entre la organización técnica del trabajo y la responsabilidad ciudadana.
    \item Se inició el uso de Git y GitHub como herramientas de registro y publicación.
\end{itemize}

\section{Referencias}

\begin{itemize}
    \item Documentación oficial de Git: \url{https://git-scm.com/doc}
    \item Documentación de GitHub: \url{https://docs.github.com}
    \item Libro Pro Git: \url{https://git-scm.com/book/es/v2}
\end{itemize}

\end{document}
